\section{Experiments}
\label{sec:experiments}

We evaluate ACR-VLA on two LIBERO task suites to test three claims:
(i)~per-DOF conformal conditioning improves task success over scalar and no-uncertainty baselines;
(ii)~each component (CGS features, ACI+COC calibration, CRM architecture) contributes to the overall gain;
(iii)~the improvement concentrates on tasks where the base VLA exhibits anisotropic uncertainty.


\subsection{Experimental Setup}
\label{sec:exp_setup}

\paragraph{Benchmarks.}
We use two suites from the LIBERO benchmark~\cite{liu2023libero}:
\textbf{LIBERO-Spatial} (10 tasks requiring spatial reasoning about object placement) and \textbf{LIBERO-Object} (10 tasks requiring generalization across object instances).
Each task is evaluated over 10 rollouts with a maximum episode length of 300 steps, yielding 100 episodes per suite.
Success is binary: an episode succeeds if the task predicate is satisfied before the step limit.

\paragraph{Base VLA.}
We build on SmolVLA~\cite{shukor2025smolvla}, a 450M-parameter flow-matching VLA that predicts action chunks of length 50.
We use the community-fine-tuned checkpoint \texttt{lerobot/smolvla\_libero}.
The policy receives dual-camera RGB input (agentview + eye-in-hand, $256{\times}256$) and 8D proprioceptive state.
Actions are 7D: 6D end-effector pose (3 translation + 3 rotation) plus 1D gripper.
The gripper dimension is handled outside the conformal pipeline via direct thresholding.

\paragraph{CGS configuration.}
At each control step, we draw $K{=}10$ action chunks via independent flow-matching rollouts and retain the first $H{=}5$ steps of each chunk (Section~\ref{sec:cgs}).
The Ledoit--Wolf shrinkage covariance~\cite{ledoit2004well} is computed per horizon step, yielding the 40D CGS feature: 30D per-DOF standard deviations ($d{\times}H = 6{\times}5$), 5D conformal radii, and 5D log-volume terms.

\paragraph{CRM architecture and training.}
The Conformal Refinement Module is a 2-layer MLP with hidden dimension 256 (${\sim}$337K parameters).
Its input concatenates the frozen observation embedding (${\sim}$960D), the 40D CGS feature, and the 30D mean action, totaling ${\sim}$1{,}030D.
The output is a 30D residual correction ($d{\times}H$).
We train CRM on 26{,}737 precomputed samples (12{,}167 from LIBERO-Spatial + 14{,}570 from LIBERO-Object), constructed from 10 demonstrations per task across all 20 tasks.
For each demonstration step, we run $K{=}10$ SmolVLA forward passes and extract CGS features.
Training data is augmented with synthetic perturbations: camera viewpoint shifts, sensor noise, and robot-state perturbation.
Optimization uses Adam with learning rate $10^{-3}$, batch size 32, for 50 epochs.
All experiments use seed 42.

\paragraph{ACI+COC configuration.}
The target miscoverage rate is $\alpha_{\mathrm{target}}{=}0.1$.
The ACI learning rate $\gamma$ and exponential decay $\rho$ for within-episode weighting are selected via grid search on a held-out validation set (details in Appendix~\ref{app:hyperparams}).

\paragraph{Baselines.}
We compare against:
\begin{itemize}[nosep]
  \item \textbf{SmolVLA (base)}: the frozen base VLA with no correction, executing the mean of $K{=}10$ samples.
  \item \textbf{Argmin selection}: draws $K{=}10$ samples and selects the one with the lowest Mahalanobis nonconformity score under Ledoit--Wolf shrinkage covariance~\cite{braun2025cgs}.
\end{itemize}
Both baselines use the same SmolVLA checkpoint and sampling budget ($K{=}10$) as ACR-VLA, ensuring that differences reflect the correction mechanism rather than the base policy or compute budget.

\paragraph{Metrics.}
The primary metric is task success rate (\%), reported as mean $\pm$ standard deviation across rollouts.
We additionally report per-task breakdowns to identify where ACR-VLA provides the largest gains.


\subsection{Main Results}
\label{sec:main_results}

Table~\ref{tab:main} compares ACR-VLA against baselines on both LIBERO suites.

\begin{table}[t]
\centering
\caption{Task success rate (\%~$\uparrow$) on LIBERO-Spatial and LIBERO-Object (10 tasks $\times$ 10 rollouts each). ACR-VLA attaches to the same frozen SmolVLA checkpoint used by all methods.}
\label{tab:main}
\small
\begin{tabular}{lcc}
\toprule
Method & LIBERO-Spatial & LIBERO-Object \\
\midrule
SmolVLA (base) & \textbf{--} & \textbf{--} \\
ReconVLA~\cite{chen2026reconvla} & \textbf{--} & \textbf{--} \\
\textbf{ACR-VLA (ours)} & \textbf{--} & \textbf{--} \\
\bottomrule
\end{tabular}
\end{table}

ACR-VLA achieves \textbf{XX.X}\% on LIBERO-Spatial and \textbf{XX.X}\% on LIBERO-Object, improving over the base SmolVLA by \textbf{XX.X}\,pp and \textbf{XX.X}\,pp respectively.
Compared to argmin selection, which uses the same $K{=}10$ sampling budget, observation-conditioned CRM provides +7.0\,pp improvement on LIBERO-Object (19.0\% vs.\ 12.0\%).
This gap confirms that learned correction can synthesize actions beyond the support of the original samples, breaking through the quality ceiling of sample selection.


\subsection{Ablation Studies}
\label{sec:ablations}

We isolate the contribution of each design choice through targeted ablations (Table~\ref{tab:ablation}).

\begin{table}[t]
\centering
\caption{Ablation results on LIBERO-Spatial and LIBERO-Object. Each row modifies one component of ACR-VLA. Success rate (\%~$\uparrow$).}
\label{tab:ablation}
\small
\begin{tabular}{llcc}
\toprule
ID & Variant & Spatial & Object \\
\midrule
-- & ACR-VLA (full) & \textbf{--} & \textbf{--} \\
\midrule
A3 & CRM w/o conformal features & \textbf{--} & \textbf{--} \\
A4 & CRM w/ scalar uncertainty & \textbf{--} & \textbf{--} \\
A5 & Per-DOF Bonferroni (no joint Mahalanobis) & \textbf{--} & \textbf{--} \\
A7 & DiU single-pass~\cite{xing2026diu} replaces CGS & \textbf{--} & \textbf{--} \\
\bottomrule
\end{tabular}
\end{table}

\paragraph{A3: CRM without conformal features.}
Removing all CGS features and conditioning CRM only on the observation embedding and mean action reduces success by \textbf{XX.X}\,pp, confirming that uncertainty conditioning is essential for effective correction rather than a redundant input.

\paragraph{A4: Scalar vs.\ per-DOF uncertainty.}
Replacing the 40D CGS feature with a single scalar (average standard deviation across all DOFs and horizon steps) reduces success by \textbf{XX.X}\,pp relative to the full model.
This mirrors the finding from the motivating experiment (Section~\ref{sec:motivating}): scalar conditioning cannot distinguish high- from low-uncertainty dimensions, leading to uniform corrections that degrade already-accurate action components.

\paragraph{A5: Per-DOF Bonferroni vs.\ joint Mahalanobis.}
Treating each of the $d{\times}H{=}30$ action dimensions independently and applying Bonferroni correction yields wider prediction regions than the joint Mahalanobis formulation.
This variant reduces success by \textbf{XX.X}\,pp, suggesting that capturing cross-DOF covariance structure produces tighter, more informative conformal regions.

\paragraph{A7: CGS (multi-sample) vs.\ DiU (single-pass).}
Replacing the $K{=}10$ sample-based CGS features with single-pass divergence-based uncertainty from DiU~\cite{xing2026diu} reduces success by \textbf{XX.X}\,pp.
DiU captures per-DOF \emph{aleatoric} uncertainty from a single forward pass, whereas CGS captures both epistemic and aleatoric components via multi-sample covariance.
The gap indicates that the epistemic component---which rises under distribution shift---is the more actionable signal for correction.

\paragraph{$K$ sensitivity.}
Figure~\ref{fig:k_sensitivity} shows ACR-VLA success as a function of the sample budget $K \in \{5, 10, 15, 20\}$.

\begin{figure}[t]
\centering
% \includegraphics[width=\linewidth]{figures/k_sensitivity.pdf}
\fbox{\parbox{0.9\linewidth}{\centering\vspace{2em}[Placeholder: $K$ sensitivity plot]\vspace{2em}}}
\caption{Task success rate vs.\ number of action samples $K$ on LIBERO-Spatial and LIBERO-Object. $K{=}10$ provides a favorable accuracy--compute trade-off.}
\label{fig:k_sensitivity}
\end{figure}

Performance improves from $K{=}5$ to $K{=}10$ as the covariance estimate becomes more reliable, then plateaus at $K{=}15$ and $K{=}20$.
This is consistent with the Ledoit--Wolf estimator's behavior: at $K{=}10$ the dimension-to-sample ratio $d/K{=}0.6$ is in the effective shrinkage regime, and additional samples yield diminishing returns.
We use $K{=}10$ as the default, as it provides the best accuracy--compute trade-off (each additional sample requires one full forward pass through the 450M-parameter base VLA).


\subsection{Analysis}
\label{sec:analysis}

\paragraph{Per-task breakdown.}
Table~\ref{tab:per_task} reports per-task success rates on LIBERO-Spatial.

\begin{table}[t]
\centering
\caption{Per-task success rate (\%~$\uparrow$) on LIBERO-Spatial (10 rollouts per task).}
\label{tab:per_task}
\small
\begin{tabular}{lccc}
\toprule
Task & SmolVLA & ReconVLA & ACR-VLA \\
\midrule
Task 1  & \textbf{--} & \textbf{--} & \textbf{--} \\
Task 2  & \textbf{--} & \textbf{--} & \textbf{--} \\
Task 3  & \textbf{--} & \textbf{--} & \textbf{--} \\
Task 4  & \textbf{--} & \textbf{--} & \textbf{--} \\
Task 5  & \textbf{--} & \textbf{--} & \textbf{--} \\
Task 6  & \textbf{--} & \textbf{--} & \textbf{--} \\
Task 7  & \textbf{--} & \textbf{--} & \textbf{--} \\
Task 8  & \textbf{--} & \textbf{--} & \textbf{--} \\
Task 9  & \textbf{--} & \textbf{--} & \textbf{--} \\
Task 10 & \textbf{--} & \textbf{--} & \textbf{--} \\
\midrule
Average & \textbf{--} & \textbf{--} & \textbf{--} \\
\bottomrule
\end{tabular}
\end{table}

ACR-VLA's gains are non-uniform across tasks.
We hypothesize that tasks where the base VLA exhibits high anisotropic uncertainty---i.e., large variance along some DOFs but not others---benefit most from per-DOF conditioning, as the CRM can selectively correct the uncertain dimensions.
Tasks where uncertainty is isotropically low (the base VLA is already confident and accurate) see smaller improvements, since there is less room for correction.

\paragraph{CGS feature visualization.}

\begin{figure}[t]
\centering
% \includegraphics[width=\linewidth]{figures/cgs_features.pdf}
\fbox{\parbox{0.9\linewidth}{\centering\vspace{2em}[Placeholder: CGS feature visualization]\vspace{2em}}}
\caption{Per-DOF standard deviations from CGS across a representative episode. High-uncertainty dimensions (e.g., $x$-translation near contact) receive larger CRM corrections, while low-uncertainty dimensions remain largely unmodified.}
\label{fig:cgs_vis}
\end{figure}

Figure~\ref{fig:cgs_vis} visualizes the per-DOF standard deviations from CGS across a representative LIBERO-Spatial episode alongside the CRM correction magnitudes.
The CRM learns to modulate its correction in proportion to the per-DOF uncertainty: dimensions with elevated $\sigma_{j,h}$ receive larger residuals, while dimensions with low variance are left approximately unchanged.
This dimension-selective behavior is the mechanism by which ACR-VLA avoids degrading already-accurate action components---a failure mode of both scalar-conditioned and unconditional correction approaches.

\paragraph{When does ACR-VLA help most?}
Across both suites, the largest improvements occur in episodes where the base VLA produces actions with high translational variance but low rotational variance (or vice versa).
In these anisotropic-uncertainty regimes, scalar methods like ReconVLA cannot distinguish the uncertain from the confident dimensions and either under-correct or over-correct uniformly.
ACR-VLA's per-DOF conditioning resolves this ambiguity, directing correction effort where it is needed.
Conversely, in episodes where the base VLA is uniformly confident, ACR-VLA's corrections are small and success rates match the baseline, indicating that the module does not introduce harmful corrections when the base policy is already performing well.
